% Part: second-order-logic
% Chapter: metatheory
% Section: compactness

\documentclass[../../../include/open-logic-section]{subfiles}

\begin{document}

\olfileid[pt]{sol}{met}{lst}

\olsection{O teorema de L\"owenheim--Skolem falha para a lógica de segunda ordem}

\begin{explain}
O teorema (descendente) de L\"owenheim--Skolem afirma que todo conjunto de
!!{sentence}s que tenha um modelo infinito tem !!a{enumerable} modelo. Ele
também é consequência do teorema da completude: a prova de completude gera
um modelo para qualquer conjunto consistente de !!{sentence}s, e esse
modelo é !!{enumerable}. Há também um teorema ascendente de
L\"owenheim--Skolem, que garante que, se um conjunto de !!{sentence}s tem
!!a{denumerable} modelo, então também tem !!a{nonenumerable} modelo. Ambos
os teoremas falham na lógica de segunda ordem.
\end{explain}


\begin{thm}
\ollabel{thm:sol-no-ls} O teorema de L\"owenheim--Skolem falha para a
lógica de segunda ordem: há !!{sentence}s com modelos infinitos, mas sem
modelos !!{enumerable}.
\end{thm}

\begin{proof}
Lembre-se de que
\[
\fn{Count} \ident \lexists[z][\lexists[u][\lforall[X][((X(z) \land
      \lforall[x][(X(x) \lif X(u(x)))]) \lif \lforall[x][X(x)])]]]
\]
é verdadeira em !!a{structure}~$\Struct{M}$ sse $\Domain{M}$ é
!!{enumerable}; portanto, $\lnot\fn{Count}$ é verdadeira em~$\Struct{M}$
sse $\Domain{M}$ é !!{nonenumerable}. Existem !!{structure}s desse
tipo---basta tomar qualquer conjunto !!{nonenumerable} como !!{domain}, p.ex., $\Pow{\Nat}$ ou~$\Real$. Logo, $\lnot \fn{Count}$ tem modelos
infinitos, mas não tem modelos !!{enumerable}.
\end{proof}

\begin{thm}
Há !!{sentence}s com modelos !!{denumerable}, mas sem modelos
!!{nonenumerable}.
\end{thm}

\begin{proof}
$\fn{Count} \land \fn{Inf}$ é verdadeira em $\Nat$, mas não em nenhuma
  !!{structure}~$\Struct{M}$ com $\Domain{M}$ !!{nonenumerable}.
\end{proof}


\end{document}

